\documentclass[11pt, aspectratio=169, xcolor=table,hyphens]{beamer}
\usepackage[utf8]{inputenc}
\usepackage[T1]{fontenc}
\usepackage{lmodern}
\usepackage[spanish]{babel}
\usetheme{default}
\usepackage{pdfrender}
\usepackage{tikz}
\usepackage{graphicx}
\usepackage{xcolor}
\usepackage{hyperref}
\usepackage{cite} %permite salto de línea en referencias largas
\usepackage{ragged2e}
\usepackage{smartdiagram}
\bibliographystyle{../apalike-es} %archivo de estilo apa ../apa-es en español en el directorio
\usepackage{pgfplots}
\usepackage{../template/uandes169template}

\begin{document}
	\author{Prof. Daniel Muñoz \\
		\href{mailto:dmunoz@miuandes.cl}{\texttt{dmunoz@miuandes.cl}}}
	\title{Resultados evaluación sumativa}
	\subtitle{Prueba 2}
	\date{\today} %agregar el día de hoy
	%\subject{subject}
	%\setbeamercovered{transparent}
	%\setbeamertemplate{navigation symbols}{}

	\maketitle

	\section{Resultados}

	\begin{frame}{Análisis de resultados: Edumétrico}
    	\begin{columns}
    		\begin{column}{.5\linewidth}
    			\begin{block}{Objetivos Módulo 1}
    				\begin{enumerate}
    					\justifying
    					\item Crear objetivos de aprendizaje coherentes con la taxonomía de Anderson y Krathwohl.
    					\item Crear indicadores de evaluación a partir de objetivos de aprendizaje utilizando la taxonomía de Anderson y Krathwohl.
    					\item Aplicar la tabla de especificaciones para planificar la evaluación.
    				\end{enumerate}
    			\end{block}
    		\end{column}
    		\begin{column}{.5\linewidth}
    		  \only<1>{
    			\begin{tikzpicture}
    				\begin{axis}[width=\linewidth,
    					symbolic x coords={EF1, EF2, ES1},
    					xtick=data, ylabel=Participación, bar width=20, ymin=0, ymax=100
    					]
    					\addplot[ybar,fill=blue] coordinates {
    						(EF1,   0)
    						(EF2,   3)
                (EF3,   93)
    						(ES1,   100)
    					};
    				\end{axis}
    			\end{tikzpicture}}
    		  \only<2>{
    			\begin{tikzpicture}
    				\begin{axis}[width=\linewidth,
    					symbolic x coords={OA1, OA2, OA3},
    					xtick=data, ylabel=Porcentaje de logro, bar width=20, ymin=0, ymax=100
    					]
    					\addplot[ybar,fill=blue] coordinates {
    						(OA1,   82)
    						(OA2,   87)
    						(OA3,   100)
    					};
    					\addplot[red,line legend,sharp plot,nodes near coords={},
    					update limits=false,shorten >=-3mm,shorten <=-3mm]
    					coordinates {(OA1,60) (OA3,60)}
    					node[midway,below]{exigencia 60\%};
    				\end{axis}
    			\end{tikzpicture}}
    		\end{column}
    	\end{columns}
    \end{frame}

    \begin{frame}{Lo menos logrado}
      \begin{block}{OA2-IE-Habilidad (76\%)}
        \textit{Cada IE tiene una habilidad correctamente asignada}
      \end{block}
    \end{frame}

\end{document}
